\begin{explanationbox}
\textbf{一、先从圆理解: 为什么椭圆需要补上 $\sqrt{r_1r_2}$?}

对于半径为 $R$ 的圆, 每条切线都和过切点的半径垂直,
因此\textbf{圆心到任意切线的距离始终是 $R$}.
把这个圆沿水平方向拉长、沿竖直方向缩短, 得到半长轴为 $a$、半短轴为 $b$ 的椭圆.
此时一般情况下 $OP$ 不再垂直于切线, 中心到切线的距离 $d$ 会随切点改变.
但是两个焦半径的乘积也同时改变, 正好满足
\[
 d=\frac{ab}{\sqrt{r_1r_2}}.
\]
可记成一句话: \textbf{切线离中心越远, 两条焦半径的乘积越小; 两者的变化恰好互相补偿}.
这里的 ``补偿'' 是代数恒等关系, 不是物理能量守恒.

当 $a=b=R$ (视为圆的极限情况) 时, 两焦点重合于圆心,
$r_1=r_2=R$, $d=R$, 公式退化为 $R\cdot R=R^2$.
所以 C012 可看作\textbf{圆的切线半径性质在椭圆上的推广}.
注意: 拉伸图形并不保持普通欧氏距离, 不能直接把圆的 $d=R$ 照搬给椭圆.

\textbf{二、用两个轴端点, 把公式记牢}

\begin{itemize}
\item 在长轴顶点 $P=(a,0)$: 切线为 $x=a$, 所以 $d=a$;
两焦半径是 $a+c$ 与 $a-c$, 乘积为 $a^2-c^2=b^2$.
于是 $\sqrt{r_1r_2}\,d=b\cdot a=ab$.
\item 在短轴顶点 $P=(0,b)$: 切线为 $y=b$, 所以 $d=b$;
两焦半径相等, 都是 $\sqrt{b^2+c^2}=a$.
于是 $\sqrt{r_1r_2}\,d=a\cdot b=ab$.
\end{itemize}
两端一个是 ``$b\times a$'', 另一个是 ``$a\times b$''.
中间位置虽然两项都变, 但乘积仍不变.
这两个端点适合\textbf{记忆、检验和判断最大最小值}, 单凭端点相等不能代替普遍证明.

\textbf{三、为什么是中心到切线的距离, 而不是 $OP$?}

椭圆在 $P(x_0,y_0)$ 处的切线为
\[
\frac{x_0x}{a^2}+\frac{y_0y}{b^2}=1.
\]
它的法向量可取
\[
\boldsymbol n=\left(\frac{x_0}{a^2},\frac{y_0}{b^2}\right).
\]
所谓中心到切线的距离, 是从 $O$ 向切线作垂线所得的\textbf{最短距离};
它不等于一般的 $OP=\sqrt{x_0^2+y_0^2}$.
圆的法向量平行于 $OP$, 但椭圆通常不平行, 只有轴端点等特殊位置才一致.

\textbf{四、真正直观的纯几何解释: 看焦点三角形的顶角}

设 $\varphi=\angle F_1PF_2$.
椭圆有\textbf{反射性质}: 在切点 $P$, 两条焦点连线 $PF_1,PF_2$ 与切线所成的锐角相等.
若这个角为 $\theta$, 则
\[
\theta=90^\circ-\frac\varphi2.
\]
从两个焦点分别向切线作垂线, 垂距记为 $h_1,h_2$.
由直角三角形,
\[
h_1=r_1\sin\theta,\qquad h_2=r_2\sin\theta.
\]
两个焦点都位于切线靠椭圆内部的同一侧, 而 $O$ 是焦点连线中点.
所以\textbf{中点到一条直线的有向距离是两端点有向距离的平均数}; 这里三点都在同侧, 因而
\[
d=\frac{h_1+h_2}{2}
=\frac{r_1+r_2}{2}\sin\theta
=a\cos\frac\varphi2.
\]
另一方面, 在三角形 $F_1PF_2$ 中运用余弦定理和半角公式,
\[
4r_1r_2\cos^2\frac\varphi2
=(r_1+r_2)^2-F_1F_2^2
=4a^2-4c^2=4b^2.
\]
所以 $\sqrt{r_1r_2}\cos(\varphi/2)=b$.
与 $d=a\cos(\varphi/2)$ 对照, 主公式自然出现.
\textbf{这就是 C012 的几何来源, 而不只是一次凑巧的代数化简}.

\textbf{五、与圆更接近的姊妹定理}

依旧用 $h_1,h_2$ 表示\textbf{两个焦点各自到切线的垂距}(不是 $r_1,r_2$).
由上面的几何关系,
\[
\boxed{h_1h_2=r_1r_2\sin^2\theta=b^2},
\qquad\boxed{d=\frac{h_1+h_2}{2}}.
\]
这条 ``两焦点到任意切线的垂距之积恒为 $b^2$'' 的性质非常值得关联记忆.
在圆的情形, 两焦点合并为圆心, $h_1=h_2=R$, 即 $R^2$.
椭圆中两段垂距可以不同, 但乘积始终是短半轴平方.

\textbf{六、高考题目的识别信号}

只要同时出现\textbf{椭圆切线、中心到直线的距离、切点的两个焦半径},
优先考虑 $\sqrt{r_1r_2}\,d=ab$, 不必先设斜率联立二次方程.
若给的是 $r_1+r_2=2a$ 以及 $|r_1-r_2|$, 先用平方差算 $r_1r_2$.
若出现焦点三角形顶角 $\varphi$, 则优先使用 $d=a\cos(\varphi/2)$.
\end{explanationbox}
